\documentclass[11pt,letterpaper]{article}

% Build steps (from this folder):
% 1) pdflatex Module1.tex
% 2) makeglossaries Module1
% 3) pdflatex Module1.tex
% 4) pdflatex Module1.tex

% ── Packages ──────────────────────────────────────────────────
\usepackage[margin=1in]{geometry}
\usepackage{titlesec}
\usepackage{enumitem}
\usepackage{booktabs}
\usepackage{tabularx}
\usepackage{graphicx}
\usepackage{xcolor}
\usepackage{hyperref}
\usepackage{fancyhdr}
\usepackage{amsmath}
\usepackage{tcolorbox}
\usepackage{multicol}
\usepackage{caption}
\usepackage{tikz}
\usepackage[acronym,toc,nopostdot]{glossaries}
\usetikzlibrary{shapes.geometric, arrows.meta, positioning, fit, calc}
\makenoidxglossaries%
% Shared acronym macro container for EE800/820 lesson plan modules.
% Usage: type \IDE, \FPGA, \UART, etc. First use is expanded; later uses show acronym only.
% Requires in each document preamble:
%   \usepackage[acronym,toc]{glossaries}
%   \makeglossaries
%   % Shared acronym macro container for EE800/820 lesson plan modules.
% Usage: type \IDE, \FPGA, \UART, etc. First use is expanded; later uses show acronym only.
% Requires in each document preamble:
%   \usepackage[acronym,toc]{glossaries}
%   \makeglossaries
%   % Shared acronym macro container for EE800/820 lesson plan modules.
% Usage: type \IDE, \FPGA, \UART, etc. First use is expanded; later uses show acronym only.
% Requires in each document preamble:
%   \usepackage[acronym,toc]{glossaries}
%   \makeglossaries
%   \input{../glossary}
% Print used acronyms near end of document:
%   \printglossary[type=\acronymtype,title={Acronyms Used}]

% Acronym definitions (only used entries appear in the printed glossary).
\newacronym{ask}{ASK}{Amplitude Shift Keying}
\newacronym{bga}{BGA}{Ball Grid Array}
\newacronym{bom}{BOM}{Bill of Materials}
\newacronym{bram}{BRAM}{Block Random Access Memory}
\newacronym{crc}{CRC}{Cyclic Redundancy Check}
\newacronym{css}{CSS}{Chirp Spread Spectrum}
\newacronym{dma}{DMA}{Direct Memory Access}
\newacronym{dsp}{DSP}{Digital Signal Processing}
\newacronym{eda}{EDA}{Exploratory Data Analysis}
\newacronym{fifo}{FIFO}{First In, First Out}
\newacronym{fpga}{FPGA}{Field-Programmable Gate Array}
\newacronym{fsk}{FSK}{Frequency Shift Keying}
\newacronym{fsm}{FSM}{Finite State Machine}
\newacronym{gpio}{GPIO}{General-Purpose Input/Output}
\newacronym{hal}{HAL}{Hardware Abstraction Layer}
\newacronym{i2c}{I\textsuperscript{2}C}{Inter-Integrated Circuit}
\newacronym{ide}{IDE}{Integrated Development Environment}
\newacronym{ila}{ILA}{Integrated Logic Analyzer}
\newacronym{ism}{ISM}{Industrial, Scientific, and Medical}
\newacronym{isr}{ISR}{Interrupt Service Routine}
\newacronym{lut}{LUT}{Look-Up Table}
\newacronym{mcu}{MCU}{Microcontroller Unit}
\newacronym{ml}{ML}{Machine Learning}
\newacronym{ook}{OOK}{On-Off Keying}
\newacronym{per}{PER}{Packet Error Rate}
\newacronym{rf}{RF}{Radio Frequency}
\newacronym{rssi}{RSSI}{Received Signal Strength Indicator}
\newacronym{rtl}{RTL}{Register-Transfer Level}
\newacronym{snr}{SNR}{Signal-to-Noise Ratio}
\newacronym{spi}{SPI}{Serial Peripheral Interface}
\newacronym{uart}{UART}{Universal Asynchronous Receiver/Transmitter}
\newacronym{vhdl}{VHDL}{VHSIC Hardware Description Language}

% Convenience wrappers so you can type \IDE, \RF, etc.
\providecommand{\ASK}{\gls{ask}}
\providecommand{\BGA}{\gls{bga}}
\providecommand{\BOM}{\gls{bom}}
\providecommand{\BRAM}{\gls{bram}}
\providecommand{\CRC}{\gls{crc}}
\providecommand{\CSS}{\gls{css}}
\providecommand{\DMA}{\gls{dma}}
\providecommand{\DSP}{\gls{dsp}}
\providecommand{\EDA}{\gls{eda}}
\providecommand{\FIFO}{\gls{fifo}}
\providecommand{\FPGA}{\gls{fpga}}
\providecommand{\FSK}{\gls{fsk}}
\providecommand{\FSM}{\gls{fsm}}
\providecommand{\GPIO}{\gls{gpio}}
\providecommand{\HAL}{\gls{hal}}
\providecommand{\IIC}{\gls{i2c}}
\providecommand{\IDE}{\gls{ide}}
\providecommand{\ILA}{\gls{ila}}
\providecommand{\ISM}{\gls{ism}}
\providecommand{\ISR}{\gls{isr}}
\providecommand{\LUT}{\gls{lut}}
\providecommand{\MCU}{\gls{mcu}}
\providecommand{\ML}{\gls{ml}}
\providecommand{\OOK}{\gls{ook}}
\providecommand{\PER}{\gls{per}}
\providecommand{\RF}{\gls{rf}}
\providecommand{\RSSI}{\gls{rssi}}
\providecommand{\RTL}{\gls{rtl}}
\providecommand{\SNR}{\gls{snr}}
\providecommand{\SPI}{\gls{spi}}
\providecommand{\UART}{\gls{uart}}
\providecommand{\VHDL}{\gls{vhdl}}


% Print used acronyms near end of document:
%   \printglossary[type=\acronymtype,title={Acronyms Used}]

% Acronym definitions (only used entries appear in the printed glossary).
\newacronym{ask}{ASK}{Amplitude Shift Keying}
\newacronym{bga}{BGA}{Ball Grid Array}
\newacronym{bom}{BOM}{Bill of Materials}
\newacronym{bram}{BRAM}{Block Random Access Memory}
\newacronym{crc}{CRC}{Cyclic Redundancy Check}
\newacronym{css}{CSS}{Chirp Spread Spectrum}
\newacronym{dma}{DMA}{Direct Memory Access}
\newacronym{dsp}{DSP}{Digital Signal Processing}
\newacronym{eda}{EDA}{Exploratory Data Analysis}
\newacronym{fifo}{FIFO}{First In, First Out}
\newacronym{fpga}{FPGA}{Field-Programmable Gate Array}
\newacronym{fsk}{FSK}{Frequency Shift Keying}
\newacronym{fsm}{FSM}{Finite State Machine}
\newacronym{gpio}{GPIO}{General-Purpose Input/Output}
\newacronym{hal}{HAL}{Hardware Abstraction Layer}
\newacronym{i2c}{I\textsuperscript{2}C}{Inter-Integrated Circuit}
\newacronym{ide}{IDE}{Integrated Development Environment}
\newacronym{ila}{ILA}{Integrated Logic Analyzer}
\newacronym{ism}{ISM}{Industrial, Scientific, and Medical}
\newacronym{isr}{ISR}{Interrupt Service Routine}
\newacronym{lut}{LUT}{Look-Up Table}
\newacronym{mcu}{MCU}{Microcontroller Unit}
\newacronym{ml}{ML}{Machine Learning}
\newacronym{ook}{OOK}{On-Off Keying}
\newacronym{per}{PER}{Packet Error Rate}
\newacronym{rf}{RF}{Radio Frequency}
\newacronym{rssi}{RSSI}{Received Signal Strength Indicator}
\newacronym{rtl}{RTL}{Register-Transfer Level}
\newacronym{snr}{SNR}{Signal-to-Noise Ratio}
\newacronym{spi}{SPI}{Serial Peripheral Interface}
\newacronym{uart}{UART}{Universal Asynchronous Receiver/Transmitter}
\newacronym{vhdl}{VHDL}{VHSIC Hardware Description Language}

% Convenience wrappers so you can type \IDE, \RF, etc.
\providecommand{\ASK}{\gls{ask}}
\providecommand{\BGA}{\gls{bga}}
\providecommand{\BOM}{\gls{bom}}
\providecommand{\BRAM}{\gls{bram}}
\providecommand{\CRC}{\gls{crc}}
\providecommand{\CSS}{\gls{css}}
\providecommand{\DMA}{\gls{dma}}
\providecommand{\DSP}{\gls{dsp}}
\providecommand{\EDA}{\gls{eda}}
\providecommand{\FIFO}{\gls{fifo}}
\providecommand{\FPGA}{\gls{fpga}}
\providecommand{\FSK}{\gls{fsk}}
\providecommand{\FSM}{\gls{fsm}}
\providecommand{\GPIO}{\gls{gpio}}
\providecommand{\HAL}{\gls{hal}}
\providecommand{\IIC}{\gls{i2c}}
\providecommand{\IDE}{\gls{ide}}
\providecommand{\ILA}{\gls{ila}}
\providecommand{\ISM}{\gls{ism}}
\providecommand{\ISR}{\gls{isr}}
\providecommand{\LUT}{\gls{lut}}
\providecommand{\MCU}{\gls{mcu}}
\providecommand{\ML}{\gls{ml}}
\providecommand{\OOK}{\gls{ook}}
\providecommand{\PER}{\gls{per}}
\providecommand{\RF}{\gls{rf}}
\providecommand{\RSSI}{\gls{rssi}}
\providecommand{\RTL}{\gls{rtl}}
\providecommand{\SNR}{\gls{snr}}
\providecommand{\SPI}{\gls{spi}}
\providecommand{\UART}{\gls{uart}}
\providecommand{\VHDL}{\gls{vhdl}}


% Print used acronyms near end of document:
%   \printglossary[type=\acronymtype,title={Acronyms Used}]

% Acronym definitions (only used entries appear in the printed glossary).
\newacronym{ask}{ASK}{Amplitude Shift Keying}
\newacronym{bga}{BGA}{Ball Grid Array}
\newacronym{bom}{BOM}{Bill of Materials}
\newacronym{bram}{BRAM}{Block Random Access Memory}
\newacronym{crc}{CRC}{Cyclic Redundancy Check}
\newacronym{css}{CSS}{Chirp Spread Spectrum}
\newacronym{dma}{DMA}{Direct Memory Access}
\newacronym{dsp}{DSP}{Digital Signal Processing}
\newacronym{eda}{EDA}{Exploratory Data Analysis}
\newacronym{fifo}{FIFO}{First In, First Out}
\newacronym{fpga}{FPGA}{Field-Programmable Gate Array}
\newacronym{fsk}{FSK}{Frequency Shift Keying}
\newacronym{fsm}{FSM}{Finite State Machine}
\newacronym{gpio}{GPIO}{General-Purpose Input/Output}
\newacronym{hal}{HAL}{Hardware Abstraction Layer}
\newacronym{i2c}{I\textsuperscript{2}C}{Inter-Integrated Circuit}
\newacronym{ide}{IDE}{Integrated Development Environment}
\newacronym{ila}{ILA}{Integrated Logic Analyzer}
\newacronym{ism}{ISM}{Industrial, Scientific, and Medical}
\newacronym{isr}{ISR}{Interrupt Service Routine}
\newacronym{lut}{LUT}{Look-Up Table}
\newacronym{mcu}{MCU}{Microcontroller Unit}
\newacronym{ml}{ML}{Machine Learning}
\newacronym{ook}{OOK}{On-Off Keying}
\newacronym{per}{PER}{Packet Error Rate}
\newacronym{rf}{RF}{Radio Frequency}
\newacronym{rssi}{RSSI}{Received Signal Strength Indicator}
\newacronym{rtl}{RTL}{Register-Transfer Level}
\newacronym{snr}{SNR}{Signal-to-Noise Ratio}
\newacronym{spi}{SPI}{Serial Peripheral Interface}
\newacronym{uart}{UART}{Universal Asynchronous Receiver/Transmitter}
\newacronym{vhdl}{VHDL}{VHSIC Hardware Description Language}

% Convenience wrappers so you can type \IDE, \RF, etc.
\providecommand{\ASK}{\gls{ask}}
\providecommand{\BGA}{\gls{bga}}
\providecommand{\BOM}{\gls{bom}}
\providecommand{\BRAM}{\gls{bram}}
\providecommand{\CRC}{\gls{crc}}
\providecommand{\CSS}{\gls{css}}
\providecommand{\DMA}{\gls{dma}}
\providecommand{\DSP}{\gls{dsp}}
\providecommand{\EDA}{\gls{eda}}
\providecommand{\FIFO}{\gls{fifo}}
\providecommand{\FPGA}{\gls{fpga}}
\providecommand{\FSK}{\gls{fsk}}
\providecommand{\FSM}{\gls{fsm}}
\providecommand{\GPIO}{\gls{gpio}}
\providecommand{\HAL}{\gls{hal}}
\providecommand{\IIC}{\gls{i2c}}
\providecommand{\IDE}{\gls{ide}}
\providecommand{\ILA}{\gls{ila}}
\providecommand{\ISM}{\gls{ism}}
\providecommand{\ISR}{\gls{isr}}
\providecommand{\LUT}{\gls{lut}}
\providecommand{\MCU}{\gls{mcu}}
\providecommand{\ML}{\gls{ml}}
\providecommand{\OOK}{\gls{ook}}
\providecommand{\PER}{\gls{per}}
\providecommand{\RF}{\gls{rf}}
\providecommand{\RSSI}{\gls{rssi}}
\providecommand{\RTL}{\gls{rtl}}
\providecommand{\SNR}{\gls{snr}}
\providecommand{\SPI}{\gls{spi}}
\providecommand{\UART}{\gls{uart}}
\providecommand{\VHDL}{\gls{vhdl}}



% ── Colors ────────────────────────────────────────────────────
\definecolor{stevens}{RGB}{163,31,52}       % Stevens Red
\definecolor{darkgray}{RGB}{60,60,60}
\definecolor{lightgray}{RGB}{240,240,240}
\definecolor{accentblue}{RGB}{30,80,140}

% ── Hyperref Setup ────────────────────────────────────────────
\hypersetup{
	colorlinks=true,
	linkcolor=stevens,
	urlcolor=accentblue,
	citecolor=darkgray,
	plainpages=false
}

% ── Header / Footer ──────────────────────────────────────────
\pagestyle{fancy}
\fancyhf{}
\fancyhead[L]{\small\textcolor{darkgray}{Stevens Institute of Technology}}
\fancyhead[R]{\small\textcolor{darkgray}{Spring 2026}}
\fancyfoot[C]{\thepage}
\renewcommand{\headrulewidth}{0.4pt}

% ── Section Formatting ───────────────────────────────────────
\titleformat{\section}
{\Large\bfseries\color{stevens}}{\thesection}{1em}{}[\vspace{-0.5em}\rule{\textwidth}{0.4pt}]
\titleformat{\subsection}
{\large\bfseries\color{darkgray}}{\thesubsection}{1em}{}

% ── tcolorbox Styles ─────────────────────────────────────────
\tcbset{
	objective/.style={colback=lightgray, colframe=stevens, 
		fonttitle=\bfseries, title=#1, boxrule=0.5pt, arc=2pt},
	infobox/.style={colback=blue!5, colframe=accentblue, 
		fonttitle=\bfseries, title=#1, boxrule=0.5pt, arc=2pt},
	deliverable/.style={colback=green!5, colframe=green!40!black,
		fonttitle=\bfseries, title=#1, boxrule=0.5pt, arc=2pt},
	warning/.style={colback=red!5, colframe=stevens,
		fonttitle=\bfseries, title=#1, boxrule=0.5pt, arc=2pt}
}

% ── Figure placeholders ───────────────────────────────────────
\newcommand{\modulefigureplaceholder}[2]{%
	\fbox{\begin{minipage}[c][2.2in][c]{#1}
		\centering\color{darkgray}\small
			\textbf{[Figure placeholder]}\\[4pt]
			\texttt{\detokenize{#2}}
	\end{minipage}}%
}

% ══════════════════════════════════════════════════════════════
\begin{document}
	
	% ── Title Page ───────────────────────────────────────────────
	\pagenumbering{Alph}
	\begin{titlepage}
		\centering
		\vspace*{2cm}
		
		{\Huge\bfseries\color{stevens} Lesson Plan: Module~1\par}
		\vspace{0.5cm}
		{\LARGE\color{darkgray} Radio Frequency Module Selection, Hardware Bring-Up,\\[4pt]
			and Cooperative Target--Interceptor Subsystem Design\par}
		
		\vspace{1.5cm}
		\rule{0.6\textwidth}{1pt}
		\vspace{1.5cm}
		
		{\Large\textbf{An Educational Framework for\\[4pt]
				AI-Driven Signal Processing}\par}
		
		\vspace{1.5cm}
		
		\begin{tabular}{rl}
			\textbf{Student:}         & Jude Eschete \\
			\textbf{Project Advisor:} & Dr.\ Bernard Yett \\
			\textbf{Program:}         & M.S.\ Computer Engineering \\
			\textbf{Institution:}     & Stevens Institute of Technology \\
			\textbf{Date Range:}      & February 28 -- March 14, 2026 \\
		\end{tabular}
		
		\vspace{2cm}
		
		\begin{tcolorbox}[colback=lightgray, colframe=darkgray, width=0.85\textwidth, arc=2pt]
			\centering\small
			This document presents a lesson plan for the first module of the capstone project,
			covering Radio Frequency transceiver selection, hardware platform validation, and the design of the
			cooperative-target / interceptor subsystem that forms the foundation for all subsequent Field-Programmable
			Gate Array processing and Machine Learning work. Two STM32+RFM95W stations exchange DATA and ACK traffic
			on a shared 906.5\,MHz channel; a third STM32+RFM95W interceptor (``Threat'') passively captures the
			bidirectional traffic, forwards it to the FPGA, and operates a button-gated PAUSE protocol that targets
			the highest-RSSI station --- closing the intercept-and-act loop without violating FCC ISM-band rules.
		\end{tcolorbox}
		
		\vfill
		{\small\textcolor{darkgray}{JEschete@stevens.edu}}
	\end{titlepage}
	
	% ── Table of Contents ────────────────────────────────────────
	\pagenumbering{roman}
	\tableofcontents
	\newpage
	\pagenumbering{arabic}
	
	% ══════════════════════════════════════════════════════════════
	\section{Module Overview}
	% ══════════════════════════════════════════════════════════════
	
	\begin{tcolorbox}[objective={Module 1 Learning Objectives}]
		Upon completion of this module, the student (or future learner following this curriculum) will be able to:
		\begin{enumerate}[label=\textbf{LO-\arabic*}]
			\item Evaluate commercially available \RF{} transceiver modules against project-specific requirements and justify a final selection with quantitative criteria.
			\item Verify end-to-end hardware communication between the STMicroelectronics STM32 Nucleo-L476RG microcontroller and the selected \RF{} module using basic firmware.
			\item Confirm \FPGA{} toolchain readiness on the Digilent Nexys A7-100T and validate a minimal loopback design.
			\item Design a complete \RF{} cooperative-target and interceptor subsystem architecture, including bidirectional DATA/ACK traffic, transmission scheduling, packet formatting, the button-gated PAUSE protocol, and the interceptor-to-FPGA data flow.
			\item Produce a hardware requirements document and subsystem design specification that downstream modules can build upon directly.
		\end{enumerate}
	\end{tcolorbox}
	
	\subsection{Context Within the Project}
	
	This module spans the first two milestones in the project timeline. It establishes the hardware baseline and subsystem architecture upon which all subsequent modules depend. The \RF{} module trade study and hardware bring-up validate that the chosen components can communicate reliably, while the cooperative-target / interceptor subsystem design defines the architecture that the \FPGA{} register-transfer level (\RTL{}) design (Module~2) and \ML{} pipeline (Module~3) will consume.
	
	\subsection{Prerequisites}
	
	Students beginning this module should have the background listed below. For students who need to catch up on any item, a recommended primer is noted; all cited works appear in full in Section~\ref{sec:resources}.
	\begin{itemize}[leftmargin=2em]
		\item Familiarity with C/C++ embedded programming (\GPIO{}, \SPI{}, \UART{}, and \IIC{}) --- see Valvano, \textit{Embedded Systems: Real-Time Interfacing to ARM Cortex-M Microcontrollers}, for peripheral fundamentals, and the STM32L476xx Reference Manual (RM0351) for the specific peripheral registers used in this module.
		\item Basic understanding of digital modulation (\ASK{}, \FSK{}, and LoRa \CSS{} concepts) --- Sklar, \textit{Digital Communications: Fundamentals and Applications}, covers the modulation theory; Semtech Application Note AN1200.22, ``LoRa Modulation Basics,'' is the concise primer for LoRa-specific \CSS{}.
		
		\newpage
		\item Experience with \VHDL{} and the \FPGA{} synthesis flow --- Chu, \textit{\FPGA{} Prototyping by \VHDL{} Examples: Xilinx Artix-7 Edition}, is the primary recommended text; Haskell and Hanna, \textit{Digital Design Using Digilent \FPGA{} Boards: \VHDL{}\,/\,Vivado Edition}, provides a Nexys-board-specific walkthrough of the Vivado toolchain.
		\item Introductory exposure to machine learning (supervised classification) --- a formal reference is introduced in Module~4; any standard introductory text on supervised classification is sufficient background here.
		\item Access to the hardware listed in Section~\ref{sec:hardware}
	\end{itemize}
	
	\newpage
	% ══════════════════════════════════════════════════════════════
	\section{Session Schedule}
	% ══════════════════════════════════════════════════════════════

	The module spans four sessions across two project milestones. Each session has a thematic focus, concrete tasks, and a deliverable checkpoint.

	\begin{table}[ht!]
		\centering
		\caption{Module 1 session breakdown.}
		\label{tab:schedule}
		\renewcommand{\arraystretch}{1.4}
		\begin{tabularx}{\textwidth}{c l X}
			\toprule
			\textbf{Session} & \textbf{Focus Area} & \textbf{Key Activities} \\
			\midrule
			1 & \RF{} Module Trade Study &
			Define selection criteria; evaluate candidate modules (RFM95W, SX1276, Ra-02, E32); perform link budget analysis; select and procure final module. \\
			2 & Hardware Bring-Up &
			Flash STM32 with \SPI{} driver firmware; verify \RF{} module register access; transmit/receive basic packets; confirm Nexys A7 toolchain and loopback; produce hardware requirements document. \\
			3 & Cooperative Target Subsystem Design &
			Define cooperative-target transmission architecture (Target~A initiates DATA, Target~B replies with ACK); design 32-byte packet format with packed DATA/ACK/PAUSE type field; implement configurable transmission scheduling and the PAUSE listener path on STM32; validate two-target ping-pong operation. \\
			4 & Interceptor (Threat) Subsystem Design &
			Design interceptor data path from \RF{} front-end to \FPGA{} input; implement passive capture of every DATA, ACK, and PAUSE frame on the channel; implement the per-Beacon-ID RSSI table and the button-gated PAUSE scheduler that targets the highest-RSSI station; define \FPGA{} ingestion interface; validate end-to-end target-to-\FPGA{} data flow including PAUSE behavioral response on the targets. \\
			\bottomrule
		\end{tabularx}
	\end{table}
	
	\newpage
	% ══════════════════════════════════════════════════════════════
	\section{Session 1: RF Module Trade Study}
	\label{sec:tradestudy}
	% ══════════════════════════════════════════════════════════════
	
	\subsection{Objectives}
	
	\begin{itemize}[leftmargin=2em]
		\item Define quantitative and qualitative selection criteria for the \RF{} transceiver module.
		\item Evaluate at least four candidate modules against those criteria.
		\item Perform a basic link budget analysis for the benchtop operating environment.
		\item Produce a justified final selection with procurement plan.
	\end{itemize}
	
	\subsection{Selection Criteria}
	
	The following table defines the evaluation dimensions. Each criterion is weighted to reflect its importance to the project's educational and technical goals.
	
	\begin{table}[ht!]
		\centering
		\caption{\RF{} module selection criteria and weights.}
		\label{tab:criteria}
		\renewcommand{\arraystretch}{1.3}
		\begin{tabularx}{\textwidth}{l c X}
			\toprule
			\textbf{Criterion} & \textbf{Weight} & \textbf{Description} \\
			\midrule
			\SPI{} Interface Compatibility & 20\% & Must support standard \SPI{} (CPOL/CPHA modes) at voltages compatible with the STM32 Nucleo (3.3\,V logic). \\
			Frequency Band & 15\% & 915\,MHz industrial, scientific, and medical (ISM) band preferred for U.S.\ regulatory compliance. \\
			Modulation Support & 15\% & Must support LoRa (\CSS{}) and ideally \FSK{}/on-off keying (\OOK{}) for multi-modulation classification exercises. \\
			Documentation \& Community & 15\% & Availability of datasheets, application notes, open-source driver libraries, and community forums. \\
			Cost \& Availability & 10\% & Per-unit cost and lead time; must be procurable from standard distributors (Digi-Key, Mouser, Adafruit). \\
			Power Consumption & 10\% & Low enough for USB-powered benchtop operation; transmit current $< 120$\,mA. \\
			Packet Engine Features & 10\% & On-chip first in, first out (\FIFO{}) buffering, cyclic redundancy check (\CRC{}), and packet handling to simplify initial firmware development. \\
			Physical Form Factor & 5\% & Breadboard- or breakout-board-friendly; no fine-pitch ball grid array (BGA) requiring reflow. \\
			\bottomrule
		\end{tabularx}
	\end{table}
	
	\newpage
	\subsection{Candidate Modules}
	
	\begin{table}[ht!]
		\centering
		\caption{Candidate \RF{} transceiver modules.}
		\label{tab:candidates}
		\renewcommand{\arraystretch}{1.3}
		\begin{tabularx}{\textwidth}{l l l X}
			\toprule
			\textbf{Module} & \textbf{Chipset} & \textbf{Freq.} & \textbf{Notes} \\
			\midrule
			Adafruit RFM95W & SX1276 & 915\,MHz & Breadboard-ready breakout; extensive tutorials. \\
			AI-Thinker Ra-02 & SX1278 & 433\,MHz & Low cost; 433\,MHz limits U.S.\ ISM exercises. \\
			EByte E32-915T20D & SX1276 & 915\,MHz & \UART{} interface (not \SPI{}); 100\,mW TX power. \\
			HopeRF RFM96W & SX1276 & 915\,MHz & SMD module; needs custom breakout or adapter. \\
			\bottomrule
		\end{tabularx}
	\end{table}
	
	\subsection{Link Budget Analysis}
	
	For the benchtop environment (transmitter--receiver separation $\leq 10$\,m indoors), students should compute:
	
	\begin{equation}
		P_{\text{rx}} = P_{\text{tx}} + G_{\text{tx}} + G_{\text{rx}} - L_{\text{path}} - L_{\text{misc}}
		\label{eq:linkbudget}
	\end{equation}
	
	where:
	\begin{itemize}[leftmargin=2em]
		\item $P_{\text{tx}}$: Transmit power (dBm), typically $+13$ to $+20$\,dBm for SX1276-based modules.
		\item $G_{\text{tx}}, G_{\text{rx}}$: Antenna gains (dBi); assume $\approx 2$\,dBi for standard whip antennas.
		\item $L_{\text{path}}$: Free-space path loss at 915\,MHz over distance $d$:
		\begin{equation}
			L_{\text{path}} = 20\log_{10}(d) + 20\log_{10}(f) - 147.55 \quad \text{[dB]}
		\end{equation}
		\item $L_{\text{misc}}$: Miscellaneous losses (cable, mismatch, indoor multipath margin $\approx 10$\,dB).
	\end{itemize}
	
	\begin{tcolorbox}[infobox={Exercise: Link Budget Worksheet}]
		Calculate the received power for the following scenario:
		\begin{itemize}
			\item $P_{\text{tx}} = +14$\,dBm, $G_{\text{tx}} = G_{\text{rx}} = 2$\,dBi
			\item $d = 5$\,m, $f = 915$\,MHz
			\item $L_{\text{misc}} = 10$\,dB (indoor margin)
		\end{itemize}
		Verify that $P_{\text{rx}}$ exceeds the SX1276 receiver sensitivity ($-137$\,dBm at SF12, 125\,kHz BW) by a comfortable margin. Document the link margin and discuss its implications for multi-beacon scaling.
	\end{tcolorbox}
	
	\begin{tcolorbox}[deliverable={Session 1 Deliverable}]
		\textbf{\RF{} Module Trade Study Report} --- A 3--5 page document containing the weighted evaluation matrix, link budget calculations, final module selection with justification, and a procurement plan (vendor, cost, estimated delivery).
	\end{tcolorbox}
	
	\newpage
	% ══════════════════════════════════════════════════════════════
	\section{Session 2: Hardware Bring-Up and Validation}
	\label{sec:bringup}
	% ══════════════════════════════════════════════════════════════
	
	\subsection{Objectives}
	
	\begin{itemize}[leftmargin=2em]
		\item Establish \SPI{} communication between the STM32 Nucleo-L476RG and the selected \RF{} module.
		\item Transmit and receive basic test packets between two cooperative-target nodes (DATA from Target~A, ACK reply from Target~B).
		\item Validate the Nexys A7-100T \FPGA{} toolchain (Vivado) with a minimal loopback design.
		\item Document all hardware interfaces, pin assignments, and power requirements.
	\end{itemize}
	
	\subsection{STM32 Firmware Bring-Up}
	
	\subsubsection{Development Environment}\label{sec:stm32_env}

	Two supported workflows are documented. Students may choose either; the course reference environment is the VS Code flow because it unifies the STM32, serial monitor, and (later) documentation/version-control workflow in a single tool.

	\paragraph{Option A --- VS Code with the official STMicroelectronics extension stack (recommended).}
	Install the \href{https://marketplace.visualstudio.com/items?itemName=stmicroelectronics.stm32-vscode-extension}{\textit{STM32 VS Code Extension}} from the Visual Studio Code marketplace. It is the official ST-maintained port of the STM32CubeIDE toolchain and automatically installs the supporting sub-extensions listed below. No separate CubeIDE installation is required.
	\begin{itemize}[leftmargin=2em]
		\item \textbf{Umbrella extension:} \texttt{stmicroelectronics.stm32-vscode-extension} (pulls the full sub-extension set).
		\item \textbf{Project management and build:} \texttt{stm32cube-ide-project-manager}, \texttt{stm32cube-ide-build-cmake}, \texttt{stm32cube-ide-build-analyzer}, and \texttt{stm32cube-ide-bundles-manager} (auto-fetches the ARM GCC toolchain and \HAL{} packs).
		\item \textbf{IntelliSense:} \texttt{stm32cube-ide-clangd} (pre-configured for ARM cross-compilation).
		\item \textbf{Debug:} \texttt{stm32cube-ide-debug-core} with the \texttt{debug-stlink-gdbserver} back end (the Nucleo-L476RG's on-board ST-LINK/V2.1 is supported out of the box; \texttt{debug-jlink-gdbserver} and \texttt{debug-generic-gdbserver} are also installed for external probes).
		\item \textbf{Peripheral register view:} \texttt{stm32cube-ide-registers} for live register inspection during debug.
		\item \textbf{RTOS awareness:} \texttt{stm32cube-ide-rtos} (unused for the Module~1 bare-metal firmware; retained for later modules if an \mbox{RTOS} is introduced).
		\newpage
		\item \textbf{Supporting general extensions:}
		\begin{itemize}[leftmargin=1.5em, topsep=2pt, itemsep=1pt]
			\item \texttt{ms-vscode.cpptools} (+ extension pack) --- fallback C/C++ language server.
			\item \texttt{ms-vscode.cmake-tools} --- CMake integration.
			\item \texttt{marus25.cortex-debug} --- alternate Cortex-M debugger with SVD, RTT, and ITM support.
			\item \texttt{dan-c-underwood.arm} --- ARM assembly syntax highlighting.
			\item \texttt{eclipse-cdt.serial-monitor} --- UART monitor for the Nucleo's virtual COM port directly inside VS Code.
		\end{itemize}
	\end{itemize}

	With this stack installed, the full workflow --- peripheral configuration (integrated CubeMX), code generation, build, flash over the on-board ST-LINK, single-step debug, peripheral register inspection, and UART monitoring --- runs entirely inside VS Code with no additional tools required.

	\paragraph{Option B --- Standalone STM32CubeIDE.}
	The classic Eclipse-based \href{https://www.st.com/en/development-tools/stm32cubeide.html}{STM32CubeIDE} (free from ST) remains fully supported. It bundles the \HAL{} libraries and STM32CubeMX code generation in a single installer and may be preferred by students already familiar with Eclipse or those who want CubeMX as a dedicated GUI application.

	\paragraph{Optional --- STM32Cube \CLT{}.}
	The \href{https://www.st.com/en/development-tools/stm32cubeclt.html}{STM32Cube Command Line Tools} package is a single installer that places the ARM GCC toolchain (\texttt{arm-none-eabi-gcc}, \texttt{-gdb}, \texttt{-objcopy}, \ldots), CMake, Ninja, \texttt{OpenOCD}, the ST-LINK GDB server, and \texttt{STM32\_Programmer\_CLI} on the system \texttt{PATH}. It is not required for the in-lab flow --- both options above handle build, flash, and debug entirely through their own integrated tooling --- but installing it is recommended if you intend to (a) script builds or flashing from a shell or CI pipeline, (b) run \texttt{arm-none-eabi-gdb} or \texttt{OpenOCD} outside of VS Code, or (c) keep a single toolchain install that survives \IDE{} upgrades. The official STM32 VS Code extension auto-discovers a \CLT{} installation when present and prefers it over the Bundles Manager download.

	\paragraph{Common target details (all options).}
	\begin{itemize}[leftmargin=2em]
		\item \textbf{Target \MCU{}:} STM32L476RG (ARM Cortex-M4, 80\,MHz, low-power).
		\item \textbf{Debug Interface:} On-board ST-LINK/V2.1 (USB) --- no external programmer required.
		\item \textbf{Serial console:} Nucleo virtual COM port via the same ST-LINK USB connection (used for debug prints and \RSSI{}/\SNR{} logging).
	\end{itemize}
	
	\subsubsection{SPI Configuration}
	
	Using the integrated STM32CubeMX peripheral configurator (available from either the VS Code extension or standalone STM32CubeIDE --- see Section~\ref{sec:stm32_env}), configure the \SPI{} peripheral as follows:
	
	\begin{table}[ht!]
		\centering
		\caption{\SPI{} peripheral configuration for \RF{} module communication.}
		\label{tab:spi}
		\renewcommand{\arraystretch}{1.3}
		\begin{tabular}{l l l}
			\toprule
			\textbf{Parameter} & \textbf{Setting} & \textbf{Rationale} \\
			\midrule
			Mode & Full-Duplex Master & \MCU{} drives clock to \RF{} module \\
			Clock Polarity (CPOL) & 0 & SX1276 default \\
			Clock Phase (CPHA) & 0 & SX1276 default \\
			Data Size & 8 bits & Register-level access \\
			Baud Rate & $\leq 10$\,MHz & SX1276 \SPI{} max = 10\,MHz \\
			NSS & Software (\GPIO{}) & Manual chip select control \\
			\bottomrule
		\end{tabular}
	\end{table}
	
	\newpage
	\subsubsection{Validation Steps}
	
	\begin{enumerate}[leftmargin=2em]
		\item \textbf{Register Read Test:} Read the SX1276 \texttt{RegVersion} register (address \texttt{0x42}). Expected return value: \texttt{0x12}. This confirms \SPI{} wiring and clock configuration.
		\item \textbf{Register Write/Read-Back:} Write a known value to a writable register (e.g., \texttt{RegFrfMsb}), read it back, and verify.
		\item \textbf{Mode Transition:} Command the module from \texttt{SLEEP} $\rightarrow$ \texttt{STANDBY} $\rightarrow$ \texttt{TX} and verify the mode bits in \texttt{RegOpMode}.
		\item \textbf{Packet Transmission:} Configure two STM32+\RF{} nodes. Transmit a known payload from Node~A; verify reception and \CRC{} pass on Node~B. Log \RSSI{} and \SNR{} values.
	\end{enumerate}
	
	\begin{tcolorbox}[warning={Common Pitfalls}]
		\begin{itemize}
			\item \textbf{Level mismatch:} Ensure the \RF{} module operates at 3.3\,V logic. The Nucleo's Arduino-compatible headers provide 3.3\,V by default, but verify with a multimeter before connecting.
			\item \textbf{\SPI{} bus contention:} If other \SPI{} peripherals share the bus, ensure proper chip-select management.
			\item \textbf{Antenna:} Always connect an antenna (or 50\,$\Omega$ dummy load) before transmitting to avoid damage to the \RF{} front-end.
			\item \textbf{Reset sequencing:} The SX1276 requires a hardware reset pulse (active low, $\geq 100$\,$\mu$s) after power-on before \SPI{} access is reliable.
		\end{itemize}
	\end{tcolorbox}
	
	\subsection{FPGA Toolchain Validation}
	
	\subsubsection{Development Environment}
	
	\begin{itemize}[leftmargin=2em]
		\item \textbf{Tool:} AMD Vivado Design Suite (WebPACK edition, free for Artix-7).
		\item \textbf{Board:} Digilent Nexys A7-100T (XC7A100T-1CSG324C).
		\item \textbf{Constraint File:} Use the Digilent-provided \texttt{Nexys-A7-100T-Master.xdc}.
	\end{itemize}
	
	\subsubsection{Loopback Design}
	
	Implement a minimal design to validate the toolchain and board:
	
	\begin{enumerate}[leftmargin=2em]
		\item \textbf{LED Echo:} Map the 16 slide switches to the 16 LEDs (combinational pass-through). Confirms pin constraints and bitstream generation.
		\item \textbf{\UART{} Loopback:} Instantiate a simple \UART{} receiver and transmitter at 115200 baud on the USB-\UART{} bridge. Echo received characters back to the host terminal. Confirms clocking, \UART{} IP, and host communication.
		\item \textbf{Clock Verification:} Use an \ILA{} (Integrated Logic Analyzer) core to probe a free-running counter driven by the 100\,MHz oscillator. Verify the expected count rate.
	\end{enumerate}
	
	\begin{tcolorbox}[deliverable={Session 2 Deliverable}]
		\textbf{Hardware Requirements Document} --- A formal document capturing:
		\begin{enumerate}
			\item \textbf{Bill of Materials (\BOM{}):} Part numbers, quantities, unit costs, vendors, and lead times for all components.
			\item \textbf{Interface Specification:} Pin mappings for STM32$\leftrightarrow$\RF{} module (\SPI{} + \GPIO{}), planned \FPGA{}$\leftrightarrow$receiver interface, \UART{} debug connections.
			\item \textbf{Power Budget:} Measured current draw for each subsystem; confirm USB power sufficiency.
			\item \textbf{Validation Evidence:} Screenshots or logic analyzer captures of successful \SPI{} transactions, packet reception logs with \RSSI{}/\SNR{}, \FPGA{} loopback test results.
			\item \textbf{Risk Register:} Identified risks (e.g., \RF{} interference in lab environment, component lead times) with mitigation strategies.
		\end{enumerate}
	\end{tcolorbox}
	
	\newpage
	% ══════════════════════════════════════════════════════════════
	\section{Session 3: Cooperative Target Subsystem Design}
	\label{sec:beacon}
	% ══════════════════════════════════════════════════════════════

	\subsection{Objectives}

	\begin{itemize}[leftmargin=2em]
		\item Define the role and behavior of each cooperative-target node within the system.
		\item Design a packet format that carries the metadata required for downstream classification \emph{and} encodes the DATA / ACK / PAUSE packet type without modifying the 32-byte radio frame layout.
		\item Implement bidirectional firmware: Target~A initiates DATA transmissions; Target~B replies with ACK; both listen continuously for PAUSE commands addressed to their own Beacon~ID.
		\item Validate two-target ping-pong operation including the PAUSE behavioral response (target prints \texttt{"I'm Jammed"} on USART2 while held in PAUSE).
	\end{itemize}

	\subsection{Cooperative Target Architecture Overview}

	Each cooperative target consists of an STM32 Nucleo-L476RG paired with an RFM95W LoRa transceiver. Both targets are TX+RX-capable and operate as a half-duplex pair: Target~A (Beacon~ID \texttt{0x01}) initiates DATA on a configurable schedule and listens briefly for an ACK from Target~B (Beacon~ID \texttt{0x02}); Target~B normally listens, replies with ACK on every successfully received DATA, and may optionally initiate its own DATA on its own schedule. Both targets continuously monitor the channel for PAUSE commands from the interceptor (Threat, Beacon~ID \texttt{0xFF}) addressed to their own ID. The cooperative-target subsystem must scale from a single initiator (Target~A only, no responder) up to two initiators (Target~A and Target~B both initiating) plus interleaved PAUSE captures, progressively increasing the demands on the interceptor and \FPGA{}.

	\begin{tcolorbox}[infobox={Design Principle: Bidirectional but Half-Duplex}]
		The cooperative-target subsystem is intentionally bidirectional --- the third radio in the system (the interceptor / Threat) is hardware-symmetric with the targets, but the targets themselves never communicate with the Threat directly. Instead, the Threat passively captures every DATA and ACK frame on the channel and, when its user button is held, emits PAUSE commands on the same channel; targets that receive a PAUSE addressed to their own ID enter a 500\,ms transmit hold-off. Each radio is half-duplex (it cannot transmit and receive simultaneously); the firmware design must explicitly switch the RFM95W between \texttt{TX}, \texttt{STANDBY}, and \texttt{RXCONTINUOUS} modes around every operation. The half-duplex constraint is itself a teaching point: students discover that an ACK that arrives before the originator switches back into RX is silently lost, and the system-level fix is a deterministic post-TX RX-window timeout rather than a tighter ACK schedule.
	\end{tcolorbox}

	\newpage
	\subsection{Packet Format Design}

	All three packet types (DATA, ACK, PAUSE) share a single fixed-length 32-byte radio frame layout. The packet type is encoded in the low two bits of the \texttt{Modulation Code} byte (offset~6); the modulation kind --- LoRa, \FSK{}, or \OOK{} --- occupies the upper six bits. This packing keeps the FPGA ingestion pipeline (Module~2) and the forwarding-frame layout (Session~4 below) unchanged across all packet types: a downstream parser sees a 32-byte payload regardless of whether the packet is DATA, ACK, or PAUSE.

	\begin{table}[ht!]
		\centering
		\caption{Radio packet format (fixed-length, 32 bytes; same layout for DATA, ACK, PAUSE).}
		\label{tab:packet}
		\renewcommand{\arraystretch}{1.3}
		\begin{tabularx}{\textwidth}{l c c X}
			\toprule
			\textbf{Field} & \textbf{Offset} & \textbf{Size} & \textbf{Description} \\
			\midrule
			Preamble Sync & 0 & 2\,B & Fixed sync word (\texttt{0xAE5A}) for frame detection. \\
			Beacon ID (sender) & 2 & 1\,B & Sender ID. Target~A = \texttt{0x01}, Target~B = \texttt{0x02}, Threat = \texttt{0xFF}. \\
			Sequence Number & 3 & 2\,B & Monotonically increasing per-sender. For ACK, echoes the acknowledged DATA sequence. \\
			TX Power & 5 & 1\,B & Transmit power setting (dBm), encoded as unsigned integer. \\
			Modulation Code & 6 & 1\,B & Packed: bits [7:2] = modulation kind (\texttt{0x01}=LoRa, \texttt{0x02}=\FSK{}, \texttt{0x03}=\OOK{}, each shifted left by 2); bits [1:0] = packet type (\texttt{00}=DATA, \texttt{01}=ACK, \texttt{10}=PAUSE, \texttt{11}=reserved). \\
			Spreading Factor & 7 & 1\,B & LoRa SF value (7--12); 0x00 if non-LoRa modulation. \\
			Bandwidth Code & 8 & 1\,B & Encoded BW setting (e.g., 0x00=125\,kHz, 0x01=250\,kHz). \\
			Timestamp & 9 & 4\,B & STM32 system tick (ms) at time of transmission. \\
			Payload (DATA/ACK) or PAUSE args & 13 & 17\,B & DATA/ACK: pseudo-random or structured pattern. PAUSE: byte~0 = target Beacon~ID being paused; byte~1 = hold-off duration in 100\,ms units; bytes 2--16 = \texttt{0xAA} fill. \\
			\CRC{}-16 & 30 & 2\,B & \CRC{}-16/CCITT-FALSE over bytes 0--29. \\
			\bottomrule
		\end{tabularx}
	\end{table}

	\begin{tcolorbox}[infobox={Design Rationale: Why Pack the Type into the Modulation Code Byte?}]
		The \texttt{Modulation Code} byte already exists in the packet for a different reason --- it labels the modulation kind so the downstream classifier has a ground-truth label for the modulation classification task. Packing the 2-bit packet type into the low bits of the same byte avoids changing the 32-byte frame layout (no downstream consumer breaks) and creates a deliberate label-leakage trap that Module~4 students must reason about: a model trained with the raw \texttt{Modulation Code} byte as a feature column trivially learns the packet type from the low two bits, even when the task is supposed to be modulation classification only. The Module~4 leakage check (LH4-B) catches this by training a logistic regression on forbidden columns alone and requiring $\geq$95\% accuracy --- a passing leakage check forces students to mask the low two bits before using \texttt{mod\_code} as a label.

		The 17-byte payload region of DATA and ACK packets carries pseudo-random bit sequences for bit-error-rate (BER) testing or structured patterns that exercise specific demodulator behavior. PAUSE packets reuse the same payload region for two purpose-specific bytes (target ID, hold-off duration) plus \texttt{0xAA} fill, keeping the radio frame size constant.
	\end{tcolorbox}

	\newpage
	\subsection{Transmission Scheduling and Bidirectional Behavior}

	The cooperative-target subsystem's traffic pattern is determined jointly by Target~A's initiator schedule, Target~B's responder behavior, and any optional initiator activity from Target~B. The scheduling strategy must balance simplicity (suitable for an educational context) with sufficient sophistication to create realistic multi-access scenarios.

	\subsubsection{Initiator Scheduling Modes}

	Students implement at least two initiator scheduling modes:

	\begin{enumerate}[leftmargin=2em]
		\item \textbf{Fixed Interval (TDMA-like):} The initiator emits DATA on a periodic schedule.
		\begin{equation}
			t_{\text{tx}}^{(k)} = T_{\text{frame}} \cdot n + \Delta t \cdot k, \quad k = 0, 1, \ldots, N-1
			\label{eq:tdma}
		\end{equation}
		where $T_{\text{frame}}$ is the frame period, $\Delta t$ is the slot duration, $k$ is the initiator index (only meaningful when both targets initiate), and $n$ is the frame counter. This mode produces predictable, non-colliding traffic and is the starting point for initial validation.

		\item \textbf{Randomized Interval (ALOHA-like):} The initiator emits DATA after a random back-off drawn from a configurable uniform distribution $[T_{\min}, T_{\max}]$. This mode introduces collisions and partial overlaps, creating a more challenging scenario for the \FPGA{} detection and the \ML{} classifier in later modules.
	\end{enumerate}

	\subsubsection{Responder Behavior}

	Target~B's responder path is event-driven, not scheduled: every successfully received DATA addressed (implicitly, by ID context) to the cooperative-target group triggers a single ACK transmission whose \texttt{Sequence Number} field echoes the received DATA's sequence. The responder does not retransmit on missed DATA; the originator detects the missed ACK via its own RX-window timeout and proceeds.

	\subsubsection{PAUSE Listener Behavior}

	Both targets continuously monitor the channel between scheduled transmissions for PAUSE packets. On receiving a PAUSE whose payload byte~0 matches the target's own Beacon~ID, the firmware sets \texttt{paused\_until\_ms = HAL\_GetTick() + (payload[1] $\times$ 100)}. While \texttt{HAL\_GetTick()} $<$ \texttt{paused\_until\_ms}: scheduled DATA transmissions are skipped, ACK responses are still permitted (the framework treats ACK as a higher-priority response to a real DATA reception), and the firmware prints \texttt{"I'm Jammed"} on USART2 once per pause-entry (throttled, not per-tick). When \texttt{paused\_until\_ms} elapses without further PAUSE refresh, the target resumes normal initiator behavior. PAUSEs addressed to a different ID are received and discarded silently --- the firmware does not log them, and they do not affect the receiving target's behavior.

	\subsubsection{Timing Implementation}

	Transmission timing is driven by the STM32's hardware timer peripheral (\texttt{TIM2} or \texttt{TIM6}) configured in periodic mode at a 1\,ms tick. Key implementation details:
	\begin{itemize}[leftmargin=2em]
		\item Use the STM32 \HAL{} timer interrupt callback (\texttt{\HAL{}\_TIM\_PeriodElapsedCallback}) to maintain a \texttt{volatile uint32\_t ms\_count} that the main-loop \FSM{} polls.
		\item For randomized initiator mode, seed the pseudo-random number generator with a unique value per target (e.g., Beacon~ID XOR'd with an ADC noise sample).
		\item Log each transmission, ACK, and PAUSE-receive event over USART2 with a millisecond timestamp for ground-truth correlation with the interceptor's captured dataset.
	\end{itemize}

	\subsection{Cooperative Target Firmware Architecture}

	The cooperative-target firmware should be structured as an \FSM{} with explicit states for the bidirectional path:

	\begin{figure}[ht!]
		\centering
		\modulefigureplaceholder{0.85\textwidth}{Images/BFA.png}
		\caption{Cooperative-target firmware state machine. Target~A: \texttt{INIT}$\to$\texttt{CONFIG}$\to$\texttt{IDLE}$\to$\texttt{BUILD\_DATA}$\to$\texttt{TX\_DATA}$\to$\texttt{WAIT\_ACK}$\to$\texttt{LOG}$\to$\texttt{IDLE}, with a \texttt{PAUSED} branch entered from any state on PAUSE reception. Target~B: same skeleton with \texttt{RX\_LISTEN} replacing the scheduled-IDLE wait, and a \texttt{BUILD\_ACK}$\to$\texttt{TX\_ACK} branch on each successful DATA reception.}
		\label{fig:beacon_fsm}
	\end{figure}

	\subsection{Validation: Two-Target Ping-Pong with PAUSE Behavioral Response}

	\begin{tcolorbox}[infobox={Exercise: Bidirectional Ping-Pong Operation}]
		Configure Target~A and Target~B with the parameters below and verify correct operation:

		\medskip
		\begin{tabular}{l c c}
			\toprule
			\textbf{Parameter} & \textbf{Target~A} & \textbf{Target~B} \\
			\midrule
			Beacon ID            & \texttt{0x01} & \texttt{0x02} \\
			Role                 & Initiator     & Responder \\
			TX Power             & +14\,dBm      & +14\,dBm \\
			Spreading Factor     & SF7           & SF7 \\
			Initiator schedule   & Fixed, 1\,s   & --- (event-driven on DATA RX) \\
			\bottomrule
		\end{tabular}

		\medskip
		Using the interceptor (Threat) firmware as a passive sniffer, capture 60 seconds of traffic and verify:
		\begin{enumerate}
			\item Each Target~A DATA (\texttt{Beacon ID} = \texttt{0x01}, \texttt{pkt\_type} = DATA) is followed by a Target~B ACK (\texttt{Beacon ID} = \texttt{0x02}, \texttt{pkt\_type} = ACK) whose \texttt{Sequence Number} matches the DATA's.
			\item Sequence numbers are monotonically increasing on Target~A's DATA stream.
			\item No \CRC{} failures occur under the fixed-interval schedule.
			\item Round-trip time (DATA TX timestamp to ACK RX timestamp on Target~A) is bounded (typically $<$120\,ms at SF7 for a 32-byte packet).
		\end{enumerate}
	\end{tcolorbox}

	\begin{tcolorbox}[infobox={Exercise: PAUSE Behavioral Response}]
		With the ping-pong above running, hold the user button (B1) on the interceptor (Threat) board and verify:
		\begin{enumerate}
			\item Target~A's USART2 log prints \texttt{"I'm Jammed"} within one PAUSE cadence interval (100\,ms) of the button being pressed.
			\item Target~A's DATA transmissions stop while the button is held.
			\item Target~A resumes normal DATA transmissions within ($\sim$500\,ms +$T_{\mathrm{slot}}$) of the button being released.
			\item Target~B's log does not print \texttt{"I'm Jammed"} (the Threat targets the higher-RSSI station, which at equal-distance is whichever is closer to the Threat antenna; if both are equidistant, the choice may alternate --- this is a measurement, not a fault).
		\end{enumerate}
	\end{tcolorbox}

	\begin{tcolorbox}[deliverable={Session 3 Deliverable}]
		\textbf{Cooperative Target Subsystem Design Document} --- A document containing:
		\begin{enumerate}
			\item Packet format specification with field descriptions, byte layout diagram, and the packed \texttt{Modulation Code} bit-field decoded.
			\item Initiator scheduling design (both TDMA and ALOHA modes) with timing equations.
			\item Cooperative target firmware architecture: \FSM{} diagram showing the \texttt{WAIT\_ACK} branch (Target~A), the \texttt{BUILD\_ACK}$\to$\texttt{TX\_ACK} branch (Target~B), and the shared \texttt{PAUSED} branch.
			\item Two-target ping-pong validation results: 60-second packet captures showing matched DATA/ACK sequence numbers, round-trip latency CSV, and \texttt{"I'm Jammed"} response/recovery log.
		\end{enumerate}
	\end{tcolorbox}

	\newpage
	% ══════════════════════════════════════════════════════════════
	\section{Session 4: Interceptor (Threat) Subsystem Design}
	\label{sec:receiver}
	% ══════════════════════════════════════════════════════════════

	\subsection{Objectives}

	\begin{itemize}[leftmargin=2em]
		\item Design the interceptor data path from the \RF{} module through the STM32 to the \FPGA{}.
		\item Implement passive capture of every DATA, ACK, and PAUSE frame on the shared channel.
		\item Implement the per-Beacon-ID \RSSI{} observation table and the button-gated PAUSE TX scheduler that targets the highest-\RSSI{} non-stale station.
		\item Define the interface protocol between the interceptor \MCU{} and the Nexys A7 \FPGA{}.
		\item Validate end-to-end data flow from cooperative-target transmission through interceptor capture to \FPGA{} ingestion, plus the closed-loop PAUSE behavioral response on the targets.
	\end{itemize}

	\subsection{Interceptor Architecture}

	The interceptor (``Threat'') station uses a third STM32 Nucleo-L476RG paired with an RFM95W module, hardware-symmetric with the two cooperative targets but operated under a different firmware role. Unlike the targets, the interceptor handles packets from multiple asynchronous senders without ever participating in the DATA/ACK exchange. Its responsibilities are:

	\begin{enumerate}[leftmargin=2em]
		\item \textbf{Packet Reception:} Detect and demodulate every incoming DATA, ACK, and PAUSE packet using the SX1276's internal packet engine and \texttt{DIO0} interrupt line in continuous-RX mode.
		\item \textbf{Metadata Extraction:} Read hardware-reported \RSSI{}, \SNR{}, and frequency error from the SX1276 registers for each received packet.
		\item \textbf{Timestamping:} Capture the STM32 system tick at the moment the \texttt{RxDone} interrupt fires, providing an interceptor-side timestamp independent of the sender's claimed timestamp.
		\item \textbf{RSSI Observation Table:} Maintain a 256-entry per-Beacon-ID table of \texttt{\{int8\_t rssi\_dbm, uint32\_t last\_seen\_ms\}} updated in the ISR alongside the ring-buffer push. Entries older than a configurable staleness threshold (5\,s default) are excluded from highest-\RSSI{} lookups.
		\item \textbf{Forwarding:} Package the received packet plus interceptor-side metadata into a 43-byte structured frame and transmit it to the \FPGA{} over \UART{}.
		\item \textbf{Button-Gated PAUSE TX:} While the user button (B1, \texttt{PC13}, active-low) is held, every $\sim$100\,ms scan the \RSSI{} table, pick the non-stale Beacon~ID with the highest \RSSI{}, build a PAUSE packet targeting that ID with hold-off duration \texttt{0x05} (500\,ms), and transmit. The Threat's own PAUSE TX is logged on USART2 but \emph{not} forwarded to the \FPGA{}; the dataset captures only what the interceptor receives off-air, which intentionally usually does not include the Threat's own transmissions (own-transmit is shadowed by the half-duplex switch).
	\end{enumerate}

	\subsection{Interceptor-to-FPGA Interface}

	\subsubsection{Interface Selection}

	\begin{table}[ht!]
		\centering
		\caption{Interface options for interceptor--\FPGA{} communication.}
		\label{tab:interface}
		\renewcommand{\arraystretch}{1.3}
		\begin{tabularx}{\textwidth}{l c c X}
			\toprule
			\textbf{Interface} & \textbf{Throughput} & \textbf{Complexity} & \textbf{Trade-offs} \\
			\midrule
			\UART{} & $\leq 921.6$\,kbps & Low & Simple two-wire connection; sufficient for low station counts; well-supported in both STM32 \HAL{} and \VHDL{}. \\
			\SPI{} (\MCU{} as master) & $\leq 10$\,Mbps & Medium & Higher throughput for scaling; requires clock, MOSI, MISO, and CS lines; \FPGA{} acts as \SPI{} slave. \\
			\bottomrule
		\end{tabularx}
	\end{table}

	\begin{tcolorbox}[infobox={Recommendation}]
		Begin with \UART{} at 115200 baud for initial integration. At 32 bytes per radio packet wrapped into the 43-byte forwarding frame, \UART{} supports approximately 28 frames per second --- comfortably above the operating envelope of the two-target ping-pong configuration plus interleaved PAUSE captures. Transition to \SPI{} only when scaling beyond three concurrent stations or when latency requirements tighten in later modules.
	\end{tcolorbox}

	\subsubsection{Interceptor-to-FPGA Frame Format}

	The interceptor wraps every received radio packet (DATA, ACK, or PAUSE) with additional metadata before forwarding to the \FPGA{}. The forwarding-frame layout is unchanged from the original receiver design --- the FPGA pipeline (Module~2) does not need to know what the radio packet's type is, only that it received 32 bytes of packet content with valid CRC.

	\begin{table}[ht!]
		\centering
		\caption{Interceptor-to-\FPGA{} forwarding frame (43 bytes).}
		\label{tab:rxframe}
		\renewcommand{\arraystretch}{1.3}
		\begin{tabularx}{\textwidth}{l c c X}
			\toprule
			\textbf{Field} & \textbf{Offset} & \textbf{Size} & \textbf{Description} \\
			\midrule
			Start Delimiter & 0 & 1\,B & Fixed byte \texttt{0x7E} for frame synchronization. \\
			Frame Length & 1 & 1\,B & Total frame length excluding start delimiter and \CRC{}; always \texttt{0x29} (41 decimal). \\
			RX Timestamp & 2 & 4\,B & Interceptor-side system tick (ms) at \texttt{RxDone} interrupt. \\
			\RSSI{} & 6 & 1\,B & Received signal strength (dBm, two's complement). \\
			\SNR{} & 7 & 1\,B & Signal-to-noise ratio (dB, signed, 0.25\,dB resolution). \\
			Freq.\ Error & 8 & 2\,B & Frequency error estimate from SX1276 (Hz, signed 16-bit). \\
			Radio Packet & 10 & 32\,B & Complete 32-byte radio packet as received: DATA, ACK, or PAUSE per Table~\ref{tab:packet}. \\
			Frame \CRC{}-8 & 42 & 1\,B & \CRC{} over bytes 1--41 for transport integrity. \\
			\bottomrule
		\end{tabularx}
	\end{table}

	\subsection{Interceptor Firmware Architecture}

	The interceptor firmware operates in an interrupt-driven architecture with two main producers (RF \texttt{RxDone} ISR and the button-poll PAUSE scheduler) and one main consumer (the UART forwarder in the main loop):

	\begin{figure}[ht!]
		\centering
		\modulefigureplaceholder{0.9\textwidth}{Images/RFA.png}
		\caption{Interceptor firmware architecture. The \texttt{DIO0 RxDone} ISR captures packets into a ring buffer and updates the per-Beacon-ID \RSSI{} table; the main loop drains the ring buffer and forwards 43-byte frames to the \FPGA{} via UART DMA. A separate main-loop poll of \texttt{B1} drives the PAUSE TX scheduler at $\sim$100\,ms cadence whenever the button is held.}
		\label{fig:rx_fw}
	\end{figure}

	\begin{tcolorbox}[warning={Concurrency Considerations}]
		\begin{itemize}
			\item The ring buffer must be sized to handle burst arrivals. With two targets in 1\,s ping-pong (effectively two frames per second) plus PAUSE bursts at 10\,Hz while the button is held, a 16-entry ring (640\,B of beacon structs) is comfortable.
			\item Use \texttt{volatile} qualifiers and a \texttt{\_\_DMB()} memory barrier between the buffer write and the head increment to prevent reordering between the ISR producer and main-loop consumer on Cortex-M4.
			\item Do \emph{not} call \HAL{} \texttt{printf} from inside the ISR --- it blocks on UART transmit and will drop \RF{} packets. All \texttt{printf} debug output happens from the main loop after dequeue.
			\item The PAUSE TX path must serialize with the receive path: switching the radio into \texttt{TX} mode for a PAUSE blocks reception for $\sim$50\,ms (the SF7 airtime of a 32-byte packet). At 10\,Hz this consumes $\sim$50\% of the channel time --- acceptable while the button is held, but a deliberate measurement consideration.
		\end{itemize}
	\end{tcolorbox}

	\subsection{FPGA Ingestion Interface}

	On the \FPGA{} side, the interceptor data is captured by a simple \UART{} receiver module (or \SPI{} slave, if that interface is selected). The Session~4 scope is limited to:

	\begin{enumerate}[leftmargin=2em]
		\item \textbf{\UART{} Receiver IP:} Implement or instantiate a \UART{} RX module on the Nexys A7 that receives bytes at the configured baud rate and writes them to a small \FIFO{}.
		\item \textbf{Frame Detection:} Implement start-delimiter detection (\texttt{0x7E}) and frame-length extraction to identify complete frames in the byte stream.
		\item \textbf{Verification:} Display received Beacon~IDs on the Nexys A7 seven-segment display or LEDs, confirming end-to-end data flow from cooperative target through interceptor to \FPGA{} fabric.
	\end{enumerate}

	\begin{tcolorbox}[infobox={Exercise: End-to-End Data Flow Validation}]
		With Target~A and Target~B running the ping-pong configuration from Session~3:
		\begin{enumerate}
			\item Verify that the \FPGA{} receives frames from both Beacon~IDs by displaying alternating Beacon~IDs on the seven-segment display (\texttt{01} and \texttt{02} interleaving).
			\item Compare the interceptor-side \RSSI{} values displayed on the \FPGA{} with the values logged by the STM32 over USART2 debug to confirm data integrity across the interface.
			\item Hold the user button on the interceptor and verify the seven-segment display occasionally shows \texttt{FF} (the Threat's own Beacon~ID), or alternatively that the \texttt{"I'm Jammed"} log appears on the targeted Target's USART2 within 100\,ms of button press.
			\item Introduce a deliberate \CRC{} error (corrupt one byte in the interceptor firmware before forwarding) and confirm the \FPGA{}'s frame \CRC{} check rejects the corrupted frame.
		\end{enumerate}
	\end{tcolorbox}

	\subsection{System Architecture Diagram}

	Figure~\ref{fig:arch} shows the complete subsystem architecture validated by the end of this module.

	\begin{figure}[ht!]
		\centering
		\modulefigureplaceholder{\textwidth}{Images/Module1Subsystem.png}
		\caption{Module~1 subsystem architecture. Solid blocks and arrows indicate components and data paths validated by end of Session~4: bidirectional DATA/ACK between Target~A and Target~B, passive capture by the interceptor (Threat), button-gated PAUSE TX from the Threat back to the targets, and UART forwarding to the FPGA. The dashed \ML{} Pipeline block indicates a downstream module.}
		\label{fig:arch}
	\end{figure}

	\begin{tcolorbox}[deliverable={Session 4 Deliverable}]
		\textbf{Interceptor Subsystem Design Document} --- A document containing:
		\begin{enumerate}
			\item Interceptor architecture description and data flow diagram.
			\item Interceptor-to-\FPGA{} interface selection rationale and frame format specification.
			\item Interceptor firmware design: interrupt architecture, ring buffer sizing, RSSI observation table, button-gated PAUSE TX scheduler with cadence and hold-off rationale.
			\item \FPGA{} ingestion module description (\UART{} RX, frame detection, verification outputs).
			\item End-to-end validation test results: cooperative-target ping-pong $\to$ interceptor $\to$ \FPGA{} data flow confirmation, and PAUSE behavioral response measurements (response and recovery latency).
		\end{enumerate}
	\end{tcolorbox}

	\newpage
	% ══════════════════════════════════════════════════════════════
	\section{Hardware Platform Reference}
	\label{sec:hardware}
	% ══════════════════════════════════════════════════════════════

	\begin{table}[ht!]
		\centering
		\caption{Required hardware components for Module 1.}
		\label{tab:hardware}
		\renewcommand{\arraystretch}{1.3}
		\begin{tabularx}{\textwidth}{l l X}
			\toprule
			\textbf{Component} & \textbf{Qty} & \textbf{Purpose} \\
			\midrule
			STM32 Nucleo-L476RG & $\geq 3$ & Two cooperative-target MCUs (Target~A and Target~B) + one interceptor (Threat) \MCU{} \\
			\RF{} Transceiver Module & $\geq 3$ & Sub-GHz LoRa communication (RFM95W, 906.5\,MHz operating carrier in the U.S.\ 915\,MHz ISM band) \\
			915\,MHz Spring Antenna & $\geq 3$ & Required for \RF{} transmission/reception \\
			Digilent Nexys A7-100T & 1 & \FPGA{} development board (Artix-7, 100T) \\
			USB Micro-B Cables & 4--5 & Power and programming for Nucleo and Nexys boards \\
			Breadboard + Jumper Wires & --- & Prototyping connections (\SPI{}, \GPIO{}, power, \UART{}) \\
			Logic Analyzer (optional) & 1 & \SPI{}/\UART{} bus debugging and timing verification \\
			\bottomrule
		\end{tabularx}
	\end{table}

	% ══════════════════════════════════════════════════════════════
	\section{Assessment Rubric}
	% ══════════════════════════════════════════════════════════════

	\begin{table}[ht!]
		\centering
		\caption{Module 1 assessment rubric.}
		\label{tab:rubric}
		\renewcommand{\arraystretch}{1.3}
		\begin{tabularx}{\textwidth}{X c c}
			\toprule
			\textbf{Assessment Item} & \textbf{Weight} & \textbf{Due} \\
			\midrule
			\RF{} Module Trade Study Report (weighted matrix, link budget, justification) & 20\% & Session 1 \\
			Hardware Requirements Document (\BOM{}, interfaces, power budget, validation evidence) & 20\% & Session 2 \\
			Cooperative Target Subsystem Design (packet format with packed type field, ping-pong firmware, PAUSE listener, two-target tests) & 30\% & Session 3 \\
			Interceptor Subsystem Design (data path, RSSI table, PAUSE scheduler, end-to-end + PAUSE behavioral validation) & 30\% & Session 4 \\
			\bottomrule
		\end{tabularx}
	\end{table}
	
	
	\newpage
	% ══════════════════════════════════════════════════════════════
	\section{Recommended Resources}\label{sec:resources}
	% ══════════════════════════════════════════════════════════════

	\subsection{Datasheets and Technical References}
	
	\begin{enumerate}[leftmargin=2em]
		\item Semtech, ``SX1276/77/78/79 Datasheet,'' Rev.\ 7, 2020.
		\item STMicroelectronics, ``STM32L476xx Reference Manual,'' RM0351, Rev.\ 9, 2023.
		\item Xilinx/AMD, ``Artix-7 FPGAs Data Sheet: DC and AC Switching Characteristics,'' DS181, 2023.
		\item Digilent, ``Nexys A7 Reference Manual,'' Rev.\ C, 2022.
	\end{enumerate}
	
	\subsection{Application Notes}
	
	\begin{enumerate}[leftmargin=2em]
		\item Semtech AN1200.22, ``LoRa Modulation Basics.''
		\item Semtech AN1200.13, ``SX1276/7/8 LoRa Modem Design Guide.''
		\item STMicroelectronics AN4286, ``\SPI{} Protocol Used in the STM32 Bootloader.''
		\item STMicroelectronics AN4013, ``STM32 Timer Overview and Low-Power Timer Usage.''
	\end{enumerate}
	
	\subsection{Textbooks}
	
	\begin{enumerate}[leftmargin=2em]
		\item P.\ P.\ Chu, \textit{\FPGA{} Prototyping by \VHDL{} Examples: Xilinx Artix-7 Edition}, 2nd ed., Wiley, 2017.
		\item R.\ E.\ Haskell and D.\ M.\ Hanna, \textit{Digital Design Using Digilent \FPGA{} Boards: \VHDL{} / Vivado Edition}, LBE Books, 2019.
		\item C.\ \"Unsalan and B.\ Tar, \textit{Digital System Design with \FPGA{}: Implementation Using Verilog and \VHDL{}}, McGraw-Hill, 2017.
		\item B.\ Sklar, \textit{Digital Communications: Fundamentals and Applications}, 2nd ed., Prentice Hall, 2001.
		\item J.\ W.\ Valvano, \textit{Embedded Systems: Real-Time Interfacing to ARM Cortex-M Microcontrollers}, 6th ed., 2023.
	\end{enumerate}
	
	\newpage
	% ══════════════════════════════════════════════════════════════
	\section{Transition to Module 2}
	% ══════════════════════════════════════════════════════════════
	
	Upon successful completion of Module~1, students will have:
	\begin{itemize}[leftmargin=2em]
		\item A validated hardware platform with confirmed \SPI{} communication across all three nodes.
		\item A fully specified and tested cooperative-target subsystem (Target~A and Target~B) capable of bidirectional ping-pong operation with PAUSE behavioral response.
		\item An interceptor (Threat) subsystem that passively captures DATA, ACK, and PAUSE traffic, maintains a per-Beacon-ID \RSSI{} observation table, and operates the button-gated PAUSE TX scheduler.
		\item End-to-end data flow from cooperative-target transmission through interceptor to \FPGA{} frame ingestion.
		\item Comprehensive design documentation covering the 32-byte packet format with packed  \\ DATA/ACK/PAUSE type field, the 43-byte forwarding frame, the firmware architectures for both cooperative-target and interceptor roles, and the closed-loop intercept-and-act behavior of the PAUSE protocol.
	\end{itemize}

	Module~2 will build directly on these foundations with \FPGA{} \RTL{} design and signal processing simulation. The \UART{}/\SPI{} ingestion interface validated in Session~4 becomes the input to the \FPGA{} processing pipeline; the 32-byte radio packet format defines the data structures that the \RTL{} must parse; and the packed \texttt{Modulation Code} byte's low two bits become the \texttt{pkt\_type} field in the Module~2 feature vector, classified at the data-to-model interface in Module~4.
	
	\vfill
	\clearpage
	\printnoidxglossary[type=\acronymtype,title={Glossary}]

	\begin{center}
		\rule{0.5\textwidth}{0.4pt}\\[6pt]
		{\small\textcolor{darkgray}{End of Module 1 Lesson Plan}}
	\end{center}
	
\end{document}